% ============================================================
%  2.2 – Nguyên lý thiết kế hình thức  (0.5 điểm)
%  Trạng thái: ĐÃ HOÀN THIỆN
% ============================================================

\subsection{Nguyên lý nền tảng}

Thiết kế hình thức của UniJob được xây dựng trên \textbf{4 nguyên lý cốt lõi} từ lý thuyết
thiết kế giao diện hiện đại:

\begin{enumerate}
  \item \textbf{Proximity (Gần kề):} Các thông tin liên quan được nhóm lại thành
    đơn vị trực quan (card). Mỗi job card chứa đầy đủ 6--8 thông tin cốt lõi trong một
    khối bao quanh bởi viền và shadow nhẹ, giúp mắt người dùng di chuyển tự nhiên.

  \item \textbf{Alignment (Căn chỉnh):} Tất cả nội dung căn trái trên desktop, căn giữa
    trên mobile. Lưới Tailwind CSS đảm bảo căn chỉnh pixel-perfect trên mọi viewport.

  \item \textbf{Contrast (Tương phản):} Màu chủ đạo xanh lá (\texttt{\#10b981}) trên nền
    trắng đạt tỷ lệ tương phản 3.1:1 (WCAG AA cho large text). Text body \texttt{\#111827}
    trên trắng đạt 16.1:1, vượt chuẩn WCAG AAA.

  \item \textbf{Repetition (Lặp lại):} Design tokens CSS (\texttt{--color-primary},
    \texttt{--color-border}...) được dùng xuyên suốt toàn bộ codebase, đảm bảo tính nhất
    quán khi thay đổi theme mà không cần sửa từng component.
\end{enumerate}

\subsection{Hệ thống Typography}

\begin{table}[H]
\centering
\caption{Cấu trúc typography trong UniJob}
\renewcommand{\arraystretch}{1.5}
\begin{tabularx}{\textwidth}{|p{3.5cm}|p{3.5cm}|X|}
\hline
\rowcolor{emeraldhead} \th{Mức} & \th{Class Tailwind} & \th{Sử dụng} \\
\hline
Display Heading & \texttt{text-4xl font-bold} & Hero headline trang chủ \\
\hline
H1 Section      & \texttt{text-2xl font-bold} & Tiêu đề trang (JobDetail, Profile) \\
\hline
H2 Card title   & \texttt{text-lg font-semibold} & Tiêu đề job card \\
\hline
Body            & \texttt{text-sm} (14px)     & Mô tả, nội dung form \\
\hline
Caption         & \texttt{text-xs text-gray-500} & Metadata: ngày, tag, badge \\
\hline
\end{tabularx}
\end{table}

Font stack mặc định: \texttt{"Segoe UI", Roboto, Arial, sans-serif} -- không nhúng custom
font để tối ưu hiệu suất tải trang.

\subsection{Hệ thống Component}

\subsubsection{Button}

Có 3 biến thể chính:
\begin{itemize}
  \item \textbf{Primary:} nền emerald, chữ trắng -- dùng cho CTA chính (``Ứng tuyển ngay'',
    ``Tải xuống PDF'').
  \item \textbf{Outline:} viền emerald, chữ emerald -- dùng cho hành động thứ cấp
    (``Chia sẻ LinkedIn'', ``Hủy'').
  \item \textbf{Ghost:} không viền, không nền -- dùng cho điều hướng nội bộ.
\end{itemize}

Tất cả button có \texttt{disabled} state với \texttt{opacity-50} và loading spinner khi
đang xử lý bất đồng bộ.

\subsubsection{Card}

Card là đơn vị giao diện cơ bản nhất của UniJob. Mỗi card có:
\begin{itemize}
  \item Border \texttt{rounded-2xl} (16px) và shadow nhẹ \texttt{shadow-sm}.
  \item Hover state \texttt{hover:shadow-md hover:border-emerald-200} để tạo phản hồi tactile.
  \item Internal padding \texttt{p-4} hoặc \texttt{p-6} tùy context.
\end{itemize}

\subsubsection{Badge và Tag}

\begin{itemize}
  \item \textbf{Category badge:} nền xanh lá nhạt (\texttt{bg-emerald-100}), chữ đậm xanh lá.
  \item \textbf{Urgent badge:} nền cam (\texttt{bg-orange-100}), icon Zap màu cam -- thu hút
    sự chú ý cho task cần người gấp.
  \item \textbf{Skill tag:} nền xám nhạt (\texttt{bg-gray-100}), font nhỏ -- không cạnh tranh
    sự chú ý với category badge.
\end{itemize}

\subsubsection{Toast Notifications}

Sử dụng \texttt{react-hot-toast} với 3 loại thông báo:
\begin{itemize}
  \item \textbf{Success:} icon tick xanh lá -- ứng tuyển thành công, lưu profile...
  \item \textbf{Error:} icon X đỏ -- lỗi mạng, form validation fail...
  \item \textbf{Info:} icon ℹ️ -- tính năng đang phát triển (LinkedIn share).
\end{itemize}

\subsection{Micro-interactions và Animation}

\begin{itemize}
  \item \textbf{Hero fade-in:} \texttt{framer-motion} \texttt{opacity: 0 → 1, y: 20 → 0}
    trong 0.6 giây khi tải trang chủ.
  \item \textbf{Apply form slide:} form ứng tuyển hiện/ẩn mượt mà với
    \texttt{transition-all} CSS khi toggle.
  \item \textbf{Spinner khi loading:} vòng quay border-emerald trên background trắng,
    thay thế toàn bộ content khi đang fetch từ Firestore.
  \item \textbf{Skeleton loading:} \texttt{animate-pulse} cho các block gray trong khi
    dữ liệu chưa về -- tránh Layout Shift (CLS).
\end{itemize}

\subsection{Layout và Spacing System}

\begin{itemize}
  \item Spacing dùng bội số 4px: \texttt{gap-4} (16px), \texttt{gap-6} (24px),
    \texttt{p-8} (32px).
  \item \textbf{Max container width:} \texttt{max-w-7xl} (1280px) cho trang danh sách,
    \texttt{max-w-5xl} (1024px) cho trang chi tiết và form.
  \item \textbf{Grid system:} 3-col (\texttt{grid-cols-3}) cho job cards trên desktop,
    1-col trên mobile; 5-col (3:2) cho CVExport layout.
\end{itemize}
